\begin{abstract}
O dimensionamento de pontes de madeira envolve múltiplas variáveis de projeto e critérios normativos conflitantes, sendo tradicionalmente conduzido por procedimentos iterativos dependentes da experiência do projetista. Este trabalho apresenta uma abordagem de projeto inteligente para pontes de madeira roliça, na qual a modelagem paramétrica dos elementos estruturais é integralmente acoplada a uma rotina de otimização robusta multiobjetivo. O problema é formulado com duas funções objetivo antagônicas, a minimização da área total de madeira e a maximização do desempenho estrutural em serviço, sujeitas às restrições de estado limite último e de serviço prescritas pela ABNT NBR 7190:2022. A resolução emprega o Algoritmo Genético de Ordenação Não Dominada II (NSGA-II), totalmente acoplado ao modelo determinístico de verificação da longarina e do tabuleiro. A robustez é incorporada de forma nativa ao processo de busca: cada solução candidata é avaliada considerando um desvio de 5\% aplicado aos valores nominais das variáveis de projeto, e as funções objetivo e restrições são agregadas por múltiplas checagens, de modo que apenas soluções tolerantes a variações permaneçam na fronteira eficiente. A metodologia foi implementada em uma plataforma computacional de acesso livre e validada em duas pontes hipotéticas de estrada vicinal, submetidas aos trens tipo TB-240 e TB-450. Os resultados demonstram que a otimização robusta identifica a fronteira de Pareto com consumo de material significativamente inferior ao de soluções geradas por amostragem aleatória, ao mesmo tempo em que preserva margem de segurança frente às incertezas geométricas. A abordagem proposta contribui para a automação, a racionalização e a modernização do projeto de pontes de madeira.
\end{abstract}

\textbf{Palavras-chave}: pontes de madeira, projeto inteligente, otimização multiobjetivo, NSGA-II, otimização robusta, NBR 7190.
